% ============================================================
% COMMON PREAMBLE — identical block for every course (do not diverge).
% Load order in each main.tex:
%     \documentclass[11pt,a4paper]{article}
%     \usepackage{lmodern}          % house serif font
%     % ============================================================
% HOUSE STYLE — identical block for every course (do not diverge).
% Requires: xcolor, titlesec, fancyhdr, geometry, enumitem, caption
% must be loadable. accent/rulegray/rowgray colors defined here.
% Also defines the house math environments (theorem/definition/...) at the
% bottom, so a single % ============================================================
% HOUSE STYLE — identical block for every course (do not diverge).
% Requires: xcolor, titlesec, fancyhdr, geometry, enumitem, caption
% must be loadable. accent/rulegray/rowgray colors defined here.
% Also defines the house math environments (theorem/definition/...) at the
% bottom, so a single % ============================================================
% HOUSE STYLE — identical block for every course (do not diverge).
% Requires: xcolor, titlesec, fancyhdr, geometry, enumitem, caption
% must be loadable. accent/rulegray/rowgray colors defined here.
% Also defines the house math environments (theorem/definition/...) at the
% bottom, so a single \input{housestyle} covers layout + math styling.
% ============================================================

% --- Page geometry ---
\usepackage[a4paper, margin=2.5cm]{geometry}
\setlength{\parindent}{0pt}
\setlength{\parskip}{6pt plus 2pt minus 1pt}
\linespread{1.15}

% --- Colors ---
\usepackage[dvipsnames]{xcolor}
\definecolor{accent}{HTML}{1E3A5F}
\definecolor{rowgray}{HTML}{F1F5F9}
\definecolor{rulegray}{HTML}{CBD5E1}

% --- Section styling (identical everywhere) ---
\usepackage{titlesec}
\titleformat{\section}
  {\normalfont\bfseries\large\color{accent}}{\thesection}{0.5em}{}
  [\color{rulegray}\titlerule]
\titleformat{\subsection}
  {\normalfont\bfseries\normalsize\color{accent}}{\thesubsection}{0.5em}{}
\titleformat{\subsubsection}
  {\normalfont\bfseries\small\itshape\color{accent}}{\thesubsubsection}{0.5em}{}
\titlespacing{\section}{0pt}{18pt plus 4pt}{8pt}
\titlespacing{\subsection}{0pt}{12pt plus 3pt}{4pt}
\titlespacing{\subsubsection}{0pt}{8pt plus 2pt}{2pt}
\setcounter{secnumdepth}{3}
\setcounter{tocdepth}{2}

% --- Numbering within section (canonical policy for the whole collection) ---
\counterwithin{equation}{section}
\counterwithin{figure}{section}
\counterwithin{table}{section}

% --- Running heads / footer (identical everywhere) ---
\usepackage{fancyhdr}
\pagestyle{fancy}
\fancyhf{}
\fancyhead[L]{\small\nouppercase{\leftmark}}
\fancyhead[R]{\small\thepage}
\renewcommand{\headrulewidth}{0.3pt}
\renewcommand{\footrulewidth}{0pt}
\fancypagestyle{plain}{\fancyhf{}\fancyfoot[C]{\small\thepage}\renewcommand{\headrulewidth}{0pt}}

% --- Lists & captions ---
\usepackage{enumitem}
\setlist{leftmargin=1.4em, itemsep=2pt, topsep=3pt, parsep=0pt, partopsep=0pt}
\usepackage{caption}
\captionsetup{font={small}, labelfont={bf,color=accent}, labelsep=period, skip=3pt}

% --- ToC spacing ---
\usepackage{tocloft}
\setlength{\cftbeforesecskip}{0.4pt}
\setlength{\cftbeforesubsecskip}{0.2pt}
\renewcommand{\cftsecfont}{\small\bfseries}
\renewcommand{\cftsubsecfont}{\small}
\renewcommand{\cftsecpagefont}{\small\bfseries}
\renewcommand{\cftsubsecpagefont}{\small}
\setlength{\cftaftertoctitleskip}{4pt}
\renewcommand{\contentsname}{Contents}

% ============================================================
% HOUSE MATH ENVIRONMENTS (folded in from the former theorems.tex).
% Run-in accent headings, numbered within section, one shared counter so
% Definition 3.2 / Theorem 3.3 interleave. Superset of every environment the
% collection uses; a course simply ignores the ones it never calls. Placed
% after \definecolor{accent} above, which these headings need.
% ============================================================
\newcounter{housethm}[section]
\renewcommand{\thehousethm}{\thesection.\arabic{housethm}}
\makeatletter
\newenvironment{house@thm}[2]% italic body: theorem/lemma/proposition/corollary/claim
  {\par\addvspace{9pt}\refstepcounter{housethm}%
   \noindent{\bfseries\color{accent}#1~\thehousethm
     \if\relax\detokenize{#2}\relax\else\ (#2)\fi.}\hspace{0.4em}\itshape\ignorespaces}
  {\par\addvspace{9pt}}
\newenvironment{house@up}[2]% upright body: definition/notation/example/method/axiom/assumption/fact/remark
  {\par\addvspace{9pt}\refstepcounter{housethm}%
   \noindent{\bfseries\color{accent}#1~\thehousethm
     \if\relax\detokenize{#2}\relax\else\ (#2)\fi.}\hspace{0.4em}\ignorespaces}
  {\par\addvspace{9pt}}
\newenvironment{theorem}[1][]{\house@thm{Theorem}{#1}}{\endhouse@thm}
\newenvironment{lemma}[1][]{\house@thm{Lemma}{#1}}{\endhouse@thm}
\newenvironment{proposition}[1][]{\house@thm{Proposition}{#1}}{\endhouse@thm}
\newenvironment{corollary}[1][]{\house@thm{Corollary}{#1}}{\endhouse@thm}
\newenvironment{claim}[1][]{\house@thm{Claim}{#1}}{\endhouse@thm}
\newenvironment{definition}[1][]{\house@up{Definition}{#1}}{\endhouse@up}
\newenvironment{notation}[1][]{\house@up{Notation}{#1}}{\endhouse@up}
\newenvironment{example}[1][]{\house@up{Example}{#1}}{\endhouse@up}
\newenvironment{method}[1][]{\house@up{Method}{#1}}{\endhouse@up}
\newenvironment{axiom}[1][]{\house@up{Axiom}{#1}}{\endhouse@up}
\newenvironment{assumption}[1][]{\house@up{Assumption}{#1}}{\endhouse@up}
\newenvironment{fact}[1][]{\house@up{Fact}{#1}}{\endhouse@up}
\newenvironment{remark}[1][]{\house@up{Remark}{#1}}{\endhouse@up}
\makeatother
% \proofname exists only when amsthm/amsproof is loaded; define-or-override so
% this block stays self-contained in courses that load neither.
\providecommand{\proofname}{Proof}
\renewcommand{\proofname}{\color{accent}Proof}
 covers layout + math styling.
% ============================================================

% --- Page geometry ---
\usepackage[a4paper, margin=2.5cm]{geometry}
\setlength{\parindent}{0pt}
\setlength{\parskip}{6pt plus 2pt minus 1pt}
\linespread{1.15}

% --- Colors ---
\usepackage[dvipsnames]{xcolor}
\definecolor{accent}{HTML}{1E3A5F}
\definecolor{rowgray}{HTML}{F1F5F9}
\definecolor{rulegray}{HTML}{CBD5E1}

% --- Section styling (identical everywhere) ---
\usepackage{titlesec}
\titleformat{\section}
  {\normalfont\bfseries\large\color{accent}}{\thesection}{0.5em}{}
  [\color{rulegray}\titlerule]
\titleformat{\subsection}
  {\normalfont\bfseries\normalsize\color{accent}}{\thesubsection}{0.5em}{}
\titleformat{\subsubsection}
  {\normalfont\bfseries\small\itshape\color{accent}}{\thesubsubsection}{0.5em}{}
\titlespacing{\section}{0pt}{18pt plus 4pt}{8pt}
\titlespacing{\subsection}{0pt}{12pt plus 3pt}{4pt}
\titlespacing{\subsubsection}{0pt}{8pt plus 2pt}{2pt}
\setcounter{secnumdepth}{3}
\setcounter{tocdepth}{2}

% --- Numbering within section (canonical policy for the whole collection) ---
\counterwithin{equation}{section}
\counterwithin{figure}{section}
\counterwithin{table}{section}

% --- Running heads / footer (identical everywhere) ---
\usepackage{fancyhdr}
\pagestyle{fancy}
\fancyhf{}
\fancyhead[L]{\small\nouppercase{\leftmark}}
\fancyhead[R]{\small\thepage}
\renewcommand{\headrulewidth}{0.3pt}
\renewcommand{\footrulewidth}{0pt}
\fancypagestyle{plain}{\fancyhf{}\fancyfoot[C]{\small\thepage}\renewcommand{\headrulewidth}{0pt}}

% --- Lists & captions ---
\usepackage{enumitem}
\setlist{leftmargin=1.4em, itemsep=2pt, topsep=3pt, parsep=0pt, partopsep=0pt}
\usepackage{caption}
\captionsetup{font={small}, labelfont={bf,color=accent}, labelsep=period, skip=3pt}

% --- ToC spacing ---
\usepackage{tocloft}
\setlength{\cftbeforesecskip}{0.4pt}
\setlength{\cftbeforesubsecskip}{0.2pt}
\renewcommand{\cftsecfont}{\small\bfseries}
\renewcommand{\cftsubsecfont}{\small}
\renewcommand{\cftsecpagefont}{\small\bfseries}
\renewcommand{\cftsubsecpagefont}{\small}
\setlength{\cftaftertoctitleskip}{4pt}
\renewcommand{\contentsname}{Contents}

% ============================================================
% HOUSE MATH ENVIRONMENTS (folded in from the former theorems.tex).
% Run-in accent headings, numbered within section, one shared counter so
% Definition 3.2 / Theorem 3.3 interleave. Superset of every environment the
% collection uses; a course simply ignores the ones it never calls. Placed
% after \definecolor{accent} above, which these headings need.
% ============================================================
\newcounter{housethm}[section]
\renewcommand{\thehousethm}{\thesection.\arabic{housethm}}
\makeatletter
\newenvironment{house@thm}[2]% italic body: theorem/lemma/proposition/corollary/claim
  {\par\addvspace{9pt}\refstepcounter{housethm}%
   \noindent{\bfseries\color{accent}#1~\thehousethm
     \if\relax\detokenize{#2}\relax\else\ (#2)\fi.}\hspace{0.4em}\itshape\ignorespaces}
  {\par\addvspace{9pt}}
\newenvironment{house@up}[2]% upright body: definition/notation/example/method/axiom/assumption/fact/remark
  {\par\addvspace{9pt}\refstepcounter{housethm}%
   \noindent{\bfseries\color{accent}#1~\thehousethm
     \if\relax\detokenize{#2}\relax\else\ (#2)\fi.}\hspace{0.4em}\ignorespaces}
  {\par\addvspace{9pt}}
\newenvironment{theorem}[1][]{\house@thm{Theorem}{#1}}{\endhouse@thm}
\newenvironment{lemma}[1][]{\house@thm{Lemma}{#1}}{\endhouse@thm}
\newenvironment{proposition}[1][]{\house@thm{Proposition}{#1}}{\endhouse@thm}
\newenvironment{corollary}[1][]{\house@thm{Corollary}{#1}}{\endhouse@thm}
\newenvironment{claim}[1][]{\house@thm{Claim}{#1}}{\endhouse@thm}
\newenvironment{definition}[1][]{\house@up{Definition}{#1}}{\endhouse@up}
\newenvironment{notation}[1][]{\house@up{Notation}{#1}}{\endhouse@up}
\newenvironment{example}[1][]{\house@up{Example}{#1}}{\endhouse@up}
\newenvironment{method}[1][]{\house@up{Method}{#1}}{\endhouse@up}
\newenvironment{axiom}[1][]{\house@up{Axiom}{#1}}{\endhouse@up}
\newenvironment{assumption}[1][]{\house@up{Assumption}{#1}}{\endhouse@up}
\newenvironment{fact}[1][]{\house@up{Fact}{#1}}{\endhouse@up}
\newenvironment{remark}[1][]{\house@up{Remark}{#1}}{\endhouse@up}
\makeatother
% \proofname exists only when amsthm/amsproof is loaded; define-or-override so
% this block stays self-contained in courses that load neither.
\providecommand{\proofname}{Proof}
\renewcommand{\proofname}{\color{accent}Proof}
 covers layout + math styling.
% ============================================================

% --- Page geometry ---
\usepackage[a4paper, margin=2.5cm]{geometry}
\setlength{\parindent}{0pt}
\setlength{\parskip}{6pt plus 2pt minus 1pt}
\linespread{1.15}

% --- Colors ---
\usepackage[dvipsnames]{xcolor}
\definecolor{accent}{HTML}{1E3A5F}
\definecolor{rowgray}{HTML}{F1F5F9}
\definecolor{rulegray}{HTML}{CBD5E1}

% --- Section styling (identical everywhere) ---
\usepackage{titlesec}
\titleformat{\section}
  {\normalfont\bfseries\large\color{accent}}{\thesection}{0.5em}{}
  [\color{rulegray}\titlerule]
\titleformat{\subsection}
  {\normalfont\bfseries\normalsize\color{accent}}{\thesubsection}{0.5em}{}
\titleformat{\subsubsection}
  {\normalfont\bfseries\small\itshape\color{accent}}{\thesubsubsection}{0.5em}{}
\titlespacing{\section}{0pt}{18pt plus 4pt}{8pt}
\titlespacing{\subsection}{0pt}{12pt plus 3pt}{4pt}
\titlespacing{\subsubsection}{0pt}{8pt plus 2pt}{2pt}
\setcounter{secnumdepth}{3}
\setcounter{tocdepth}{2}

% --- Numbering within section (canonical policy for the whole collection) ---
\counterwithin{equation}{section}
\counterwithin{figure}{section}
\counterwithin{table}{section}

% --- Running heads / footer (identical everywhere) ---
\usepackage{fancyhdr}
\pagestyle{fancy}
\fancyhf{}
\fancyhead[L]{\small\nouppercase{\leftmark}}
\fancyhead[R]{\small\thepage}
\renewcommand{\headrulewidth}{0.3pt}
\renewcommand{\footrulewidth}{0pt}
\fancypagestyle{plain}{\fancyhf{}\fancyfoot[C]{\small\thepage}\renewcommand{\headrulewidth}{0pt}}

% --- Lists & captions ---
\usepackage{enumitem}
\setlist{leftmargin=1.4em, itemsep=2pt, topsep=3pt, parsep=0pt, partopsep=0pt}
\usepackage{caption}
\captionsetup{font={small}, labelfont={bf,color=accent}, labelsep=period, skip=3pt}

% --- ToC spacing ---
\usepackage{tocloft}
\setlength{\cftbeforesecskip}{0.4pt}
\setlength{\cftbeforesubsecskip}{0.2pt}
\renewcommand{\cftsecfont}{\small\bfseries}
\renewcommand{\cftsubsecfont}{\small}
\renewcommand{\cftsecpagefont}{\small\bfseries}
\renewcommand{\cftsubsecpagefont}{\small}
\setlength{\cftaftertoctitleskip}{4pt}
\renewcommand{\contentsname}{Contents}

% ============================================================
% HOUSE MATH ENVIRONMENTS (folded in from the former theorems.tex).
% Run-in accent headings, numbered within section, one shared counter so
% Definition 3.2 / Theorem 3.3 interleave. Superset of every environment the
% collection uses; a course simply ignores the ones it never calls. Placed
% after \definecolor{accent} above, which these headings need.
% ============================================================
\newcounter{housethm}[section]
\renewcommand{\thehousethm}{\thesection.\arabic{housethm}}
\makeatletter
\newenvironment{house@thm}[2]% italic body: theorem/lemma/proposition/corollary/claim
  {\par\addvspace{9pt}\refstepcounter{housethm}%
   \noindent{\bfseries\color{accent}#1~\thehousethm
     \if\relax\detokenize{#2}\relax\else\ (#2)\fi.}\hspace{0.4em}\itshape\ignorespaces}
  {\par\addvspace{9pt}}
\newenvironment{house@up}[2]% upright body: definition/notation/example/method/axiom/assumption/fact/remark
  {\par\addvspace{9pt}\refstepcounter{housethm}%
   \noindent{\bfseries\color{accent}#1~\thehousethm
     \if\relax\detokenize{#2}\relax\else\ (#2)\fi.}\hspace{0.4em}\ignorespaces}
  {\par\addvspace{9pt}}
\newenvironment{theorem}[1][]{\house@thm{Theorem}{#1}}{\endhouse@thm}
\newenvironment{lemma}[1][]{\house@thm{Lemma}{#1}}{\endhouse@thm}
\newenvironment{proposition}[1][]{\house@thm{Proposition}{#1}}{\endhouse@thm}
\newenvironment{corollary}[1][]{\house@thm{Corollary}{#1}}{\endhouse@thm}
\newenvironment{claim}[1][]{\house@thm{Claim}{#1}}{\endhouse@thm}
\newenvironment{definition}[1][]{\house@up{Definition}{#1}}{\endhouse@up}
\newenvironment{notation}[1][]{\house@up{Notation}{#1}}{\endhouse@up}
\newenvironment{example}[1][]{\house@up{Example}{#1}}{\endhouse@up}
\newenvironment{method}[1][]{\house@up{Method}{#1}}{\endhouse@up}
\newenvironment{axiom}[1][]{\house@up{Axiom}{#1}}{\endhouse@up}
\newenvironment{assumption}[1][]{\house@up{Assumption}{#1}}{\endhouse@up}
\newenvironment{fact}[1][]{\house@up{Fact}{#1}}{\endhouse@up}
\newenvironment{remark}[1][]{\house@up{Remark}{#1}}{\endhouse@up}
\makeatother
% \proofname exists only when amsthm/amsproof is loaded; define-or-override so
% this block stays self-contained in courses that load neither.
\providecommand{\proofname}{Proof}
\renewcommand{\proofname}{\color{accent}Proof}
            % geometry, colors, sections, ToC, heads
%     \input{theorems}              % house math environments
%     % ============================================================
% COMMON PREAMBLE — identical block for every course (do not diverge).
% Load order in each main.tex:
%     \documentclass[11pt,a4paper]{article}
%     \usepackage{lmodern}          % house serif font
%     % ============================================================
% HOUSE STYLE — identical block for every course (do not diverge).
% Requires: xcolor, titlesec, fancyhdr, geometry, enumitem, caption
% must be loadable. accent/rulegray/rowgray colors defined here.
% Also defines the house math environments (theorem/definition/...) at the
% bottom, so a single % ============================================================
% HOUSE STYLE — identical block for every course (do not diverge).
% Requires: xcolor, titlesec, fancyhdr, geometry, enumitem, caption
% must be loadable. accent/rulegray/rowgray colors defined here.
% Also defines the house math environments (theorem/definition/...) at the
% bottom, so a single \input{housestyle} covers layout + math styling.
% ============================================================

% --- Page geometry ---
\usepackage[a4paper, margin=2.5cm]{geometry}
\setlength{\parindent}{0pt}
\setlength{\parskip}{6pt plus 2pt minus 1pt}
\linespread{1.15}

% --- Colors ---
\usepackage[dvipsnames]{xcolor}
\definecolor{accent}{HTML}{1E3A5F}
\definecolor{rowgray}{HTML}{F1F5F9}
\definecolor{rulegray}{HTML}{CBD5E1}

% --- Section styling (identical everywhere) ---
\usepackage{titlesec}
\titleformat{\section}
  {\normalfont\bfseries\large\color{accent}}{\thesection}{0.5em}{}
  [\color{rulegray}\titlerule]
\titleformat{\subsection}
  {\normalfont\bfseries\normalsize\color{accent}}{\thesubsection}{0.5em}{}
\titleformat{\subsubsection}
  {\normalfont\bfseries\small\itshape\color{accent}}{\thesubsubsection}{0.5em}{}
\titlespacing{\section}{0pt}{18pt plus 4pt}{8pt}
\titlespacing{\subsection}{0pt}{12pt plus 3pt}{4pt}
\titlespacing{\subsubsection}{0pt}{8pt plus 2pt}{2pt}
\setcounter{secnumdepth}{3}
\setcounter{tocdepth}{2}

% --- Numbering within section (canonical policy for the whole collection) ---
\counterwithin{equation}{section}
\counterwithin{figure}{section}
\counterwithin{table}{section}

% --- Running heads / footer (identical everywhere) ---
\usepackage{fancyhdr}
\pagestyle{fancy}
\fancyhf{}
\fancyhead[L]{\small\nouppercase{\leftmark}}
\fancyhead[R]{\small\thepage}
\renewcommand{\headrulewidth}{0.3pt}
\renewcommand{\footrulewidth}{0pt}
\fancypagestyle{plain}{\fancyhf{}\fancyfoot[C]{\small\thepage}\renewcommand{\headrulewidth}{0pt}}

% --- Lists & captions ---
\usepackage{enumitem}
\setlist{leftmargin=1.4em, itemsep=2pt, topsep=3pt, parsep=0pt, partopsep=0pt}
\usepackage{caption}
\captionsetup{font={small}, labelfont={bf,color=accent}, labelsep=period, skip=3pt}

% --- ToC spacing ---
\usepackage{tocloft}
\setlength{\cftbeforesecskip}{0.4pt}
\setlength{\cftbeforesubsecskip}{0.2pt}
\renewcommand{\cftsecfont}{\small\bfseries}
\renewcommand{\cftsubsecfont}{\small}
\renewcommand{\cftsecpagefont}{\small\bfseries}
\renewcommand{\cftsubsecpagefont}{\small}
\setlength{\cftaftertoctitleskip}{4pt}
\renewcommand{\contentsname}{Contents}

% ============================================================
% HOUSE MATH ENVIRONMENTS (folded in from the former theorems.tex).
% Run-in accent headings, numbered within section, one shared counter so
% Definition 3.2 / Theorem 3.3 interleave. Superset of every environment the
% collection uses; a course simply ignores the ones it never calls. Placed
% after \definecolor{accent} above, which these headings need.
% ============================================================
\newcounter{housethm}[section]
\renewcommand{\thehousethm}{\thesection.\arabic{housethm}}
\makeatletter
\newenvironment{house@thm}[2]% italic body: theorem/lemma/proposition/corollary/claim
  {\par\addvspace{9pt}\refstepcounter{housethm}%
   \noindent{\bfseries\color{accent}#1~\thehousethm
     \if\relax\detokenize{#2}\relax\else\ (#2)\fi.}\hspace{0.4em}\itshape\ignorespaces}
  {\par\addvspace{9pt}}
\newenvironment{house@up}[2]% upright body: definition/notation/example/method/axiom/assumption/fact/remark
  {\par\addvspace{9pt}\refstepcounter{housethm}%
   \noindent{\bfseries\color{accent}#1~\thehousethm
     \if\relax\detokenize{#2}\relax\else\ (#2)\fi.}\hspace{0.4em}\ignorespaces}
  {\par\addvspace{9pt}}
\newenvironment{theorem}[1][]{\house@thm{Theorem}{#1}}{\endhouse@thm}
\newenvironment{lemma}[1][]{\house@thm{Lemma}{#1}}{\endhouse@thm}
\newenvironment{proposition}[1][]{\house@thm{Proposition}{#1}}{\endhouse@thm}
\newenvironment{corollary}[1][]{\house@thm{Corollary}{#1}}{\endhouse@thm}
\newenvironment{claim}[1][]{\house@thm{Claim}{#1}}{\endhouse@thm}
\newenvironment{definition}[1][]{\house@up{Definition}{#1}}{\endhouse@up}
\newenvironment{notation}[1][]{\house@up{Notation}{#1}}{\endhouse@up}
\newenvironment{example}[1][]{\house@up{Example}{#1}}{\endhouse@up}
\newenvironment{method}[1][]{\house@up{Method}{#1}}{\endhouse@up}
\newenvironment{axiom}[1][]{\house@up{Axiom}{#1}}{\endhouse@up}
\newenvironment{assumption}[1][]{\house@up{Assumption}{#1}}{\endhouse@up}
\newenvironment{fact}[1][]{\house@up{Fact}{#1}}{\endhouse@up}
\newenvironment{remark}[1][]{\house@up{Remark}{#1}}{\endhouse@up}
\makeatother
% \proofname exists only when amsthm/amsproof is loaded; define-or-override so
% this block stays self-contained in courses that load neither.
\providecommand{\proofname}{Proof}
\renewcommand{\proofname}{\color{accent}Proof}
 covers layout + math styling.
% ============================================================

% --- Page geometry ---
\usepackage[a4paper, margin=2.5cm]{geometry}
\setlength{\parindent}{0pt}
\setlength{\parskip}{6pt plus 2pt minus 1pt}
\linespread{1.15}

% --- Colors ---
\usepackage[dvipsnames]{xcolor}
\definecolor{accent}{HTML}{1E3A5F}
\definecolor{rowgray}{HTML}{F1F5F9}
\definecolor{rulegray}{HTML}{CBD5E1}

% --- Section styling (identical everywhere) ---
\usepackage{titlesec}
\titleformat{\section}
  {\normalfont\bfseries\large\color{accent}}{\thesection}{0.5em}{}
  [\color{rulegray}\titlerule]
\titleformat{\subsection}
  {\normalfont\bfseries\normalsize\color{accent}}{\thesubsection}{0.5em}{}
\titleformat{\subsubsection}
  {\normalfont\bfseries\small\itshape\color{accent}}{\thesubsubsection}{0.5em}{}
\titlespacing{\section}{0pt}{18pt plus 4pt}{8pt}
\titlespacing{\subsection}{0pt}{12pt plus 3pt}{4pt}
\titlespacing{\subsubsection}{0pt}{8pt plus 2pt}{2pt}
\setcounter{secnumdepth}{3}
\setcounter{tocdepth}{2}

% --- Numbering within section (canonical policy for the whole collection) ---
\counterwithin{equation}{section}
\counterwithin{figure}{section}
\counterwithin{table}{section}

% --- Running heads / footer (identical everywhere) ---
\usepackage{fancyhdr}
\pagestyle{fancy}
\fancyhf{}
\fancyhead[L]{\small\nouppercase{\leftmark}}
\fancyhead[R]{\small\thepage}
\renewcommand{\headrulewidth}{0.3pt}
\renewcommand{\footrulewidth}{0pt}
\fancypagestyle{plain}{\fancyhf{}\fancyfoot[C]{\small\thepage}\renewcommand{\headrulewidth}{0pt}}

% --- Lists & captions ---
\usepackage{enumitem}
\setlist{leftmargin=1.4em, itemsep=2pt, topsep=3pt, parsep=0pt, partopsep=0pt}
\usepackage{caption}
\captionsetup{font={small}, labelfont={bf,color=accent}, labelsep=period, skip=3pt}

% --- ToC spacing ---
\usepackage{tocloft}
\setlength{\cftbeforesecskip}{0.4pt}
\setlength{\cftbeforesubsecskip}{0.2pt}
\renewcommand{\cftsecfont}{\small\bfseries}
\renewcommand{\cftsubsecfont}{\small}
\renewcommand{\cftsecpagefont}{\small\bfseries}
\renewcommand{\cftsubsecpagefont}{\small}
\setlength{\cftaftertoctitleskip}{4pt}
\renewcommand{\contentsname}{Contents}

% ============================================================
% HOUSE MATH ENVIRONMENTS (folded in from the former theorems.tex).
% Run-in accent headings, numbered within section, one shared counter so
% Definition 3.2 / Theorem 3.3 interleave. Superset of every environment the
% collection uses; a course simply ignores the ones it never calls. Placed
% after \definecolor{accent} above, which these headings need.
% ============================================================
\newcounter{housethm}[section]
\renewcommand{\thehousethm}{\thesection.\arabic{housethm}}
\makeatletter
\newenvironment{house@thm}[2]% italic body: theorem/lemma/proposition/corollary/claim
  {\par\addvspace{9pt}\refstepcounter{housethm}%
   \noindent{\bfseries\color{accent}#1~\thehousethm
     \if\relax\detokenize{#2}\relax\else\ (#2)\fi.}\hspace{0.4em}\itshape\ignorespaces}
  {\par\addvspace{9pt}}
\newenvironment{house@up}[2]% upright body: definition/notation/example/method/axiom/assumption/fact/remark
  {\par\addvspace{9pt}\refstepcounter{housethm}%
   \noindent{\bfseries\color{accent}#1~\thehousethm
     \if\relax\detokenize{#2}\relax\else\ (#2)\fi.}\hspace{0.4em}\ignorespaces}
  {\par\addvspace{9pt}}
\newenvironment{theorem}[1][]{\house@thm{Theorem}{#1}}{\endhouse@thm}
\newenvironment{lemma}[1][]{\house@thm{Lemma}{#1}}{\endhouse@thm}
\newenvironment{proposition}[1][]{\house@thm{Proposition}{#1}}{\endhouse@thm}
\newenvironment{corollary}[1][]{\house@thm{Corollary}{#1}}{\endhouse@thm}
\newenvironment{claim}[1][]{\house@thm{Claim}{#1}}{\endhouse@thm}
\newenvironment{definition}[1][]{\house@up{Definition}{#1}}{\endhouse@up}
\newenvironment{notation}[1][]{\house@up{Notation}{#1}}{\endhouse@up}
\newenvironment{example}[1][]{\house@up{Example}{#1}}{\endhouse@up}
\newenvironment{method}[1][]{\house@up{Method}{#1}}{\endhouse@up}
\newenvironment{axiom}[1][]{\house@up{Axiom}{#1}}{\endhouse@up}
\newenvironment{assumption}[1][]{\house@up{Assumption}{#1}}{\endhouse@up}
\newenvironment{fact}[1][]{\house@up{Fact}{#1}}{\endhouse@up}
\newenvironment{remark}[1][]{\house@up{Remark}{#1}}{\endhouse@up}
\makeatother
% \proofname exists only when amsthm/amsproof is loaded; define-or-override so
% this block stays self-contained in courses that load neither.
\providecommand{\proofname}{Proof}
\renewcommand{\proofname}{\color{accent}Proof}
            % geometry, colors, sections, ToC, heads
%     \input{theorems}              % house math environments
%     % ============================================================
% COMMON PREAMBLE — identical block for every course (do not diverge).
% Load order in each main.tex:
%     \documentclass[11pt,a4paper]{article}
%     \usepackage{lmodern}          % house serif font
%     % ============================================================
% HOUSE STYLE — identical block for every course (do not diverge).
% Requires: xcolor, titlesec, fancyhdr, geometry, enumitem, caption
% must be loadable. accent/rulegray/rowgray colors defined here.
% Also defines the house math environments (theorem/definition/...) at the
% bottom, so a single \input{housestyle} covers layout + math styling.
% ============================================================

% --- Page geometry ---
\usepackage[a4paper, margin=2.5cm]{geometry}
\setlength{\parindent}{0pt}
\setlength{\parskip}{6pt plus 2pt minus 1pt}
\linespread{1.15}

% --- Colors ---
\usepackage[dvipsnames]{xcolor}
\definecolor{accent}{HTML}{1E3A5F}
\definecolor{rowgray}{HTML}{F1F5F9}
\definecolor{rulegray}{HTML}{CBD5E1}

% --- Section styling (identical everywhere) ---
\usepackage{titlesec}
\titleformat{\section}
  {\normalfont\bfseries\large\color{accent}}{\thesection}{0.5em}{}
  [\color{rulegray}\titlerule]
\titleformat{\subsection}
  {\normalfont\bfseries\normalsize\color{accent}}{\thesubsection}{0.5em}{}
\titleformat{\subsubsection}
  {\normalfont\bfseries\small\itshape\color{accent}}{\thesubsubsection}{0.5em}{}
\titlespacing{\section}{0pt}{18pt plus 4pt}{8pt}
\titlespacing{\subsection}{0pt}{12pt plus 3pt}{4pt}
\titlespacing{\subsubsection}{0pt}{8pt plus 2pt}{2pt}
\setcounter{secnumdepth}{3}
\setcounter{tocdepth}{2}

% --- Numbering within section (canonical policy for the whole collection) ---
\counterwithin{equation}{section}
\counterwithin{figure}{section}
\counterwithin{table}{section}

% --- Running heads / footer (identical everywhere) ---
\usepackage{fancyhdr}
\pagestyle{fancy}
\fancyhf{}
\fancyhead[L]{\small\nouppercase{\leftmark}}
\fancyhead[R]{\small\thepage}
\renewcommand{\headrulewidth}{0.3pt}
\renewcommand{\footrulewidth}{0pt}
\fancypagestyle{plain}{\fancyhf{}\fancyfoot[C]{\small\thepage}\renewcommand{\headrulewidth}{0pt}}

% --- Lists & captions ---
\usepackage{enumitem}
\setlist{leftmargin=1.4em, itemsep=2pt, topsep=3pt, parsep=0pt, partopsep=0pt}
\usepackage{caption}
\captionsetup{font={small}, labelfont={bf,color=accent}, labelsep=period, skip=3pt}

% --- ToC spacing ---
\usepackage{tocloft}
\setlength{\cftbeforesecskip}{0.4pt}
\setlength{\cftbeforesubsecskip}{0.2pt}
\renewcommand{\cftsecfont}{\small\bfseries}
\renewcommand{\cftsubsecfont}{\small}
\renewcommand{\cftsecpagefont}{\small\bfseries}
\renewcommand{\cftsubsecpagefont}{\small}
\setlength{\cftaftertoctitleskip}{4pt}
\renewcommand{\contentsname}{Contents}

% ============================================================
% HOUSE MATH ENVIRONMENTS (folded in from the former theorems.tex).
% Run-in accent headings, numbered within section, one shared counter so
% Definition 3.2 / Theorem 3.3 interleave. Superset of every environment the
% collection uses; a course simply ignores the ones it never calls. Placed
% after \definecolor{accent} above, which these headings need.
% ============================================================
\newcounter{housethm}[section]
\renewcommand{\thehousethm}{\thesection.\arabic{housethm}}
\makeatletter
\newenvironment{house@thm}[2]% italic body: theorem/lemma/proposition/corollary/claim
  {\par\addvspace{9pt}\refstepcounter{housethm}%
   \noindent{\bfseries\color{accent}#1~\thehousethm
     \if\relax\detokenize{#2}\relax\else\ (#2)\fi.}\hspace{0.4em}\itshape\ignorespaces}
  {\par\addvspace{9pt}}
\newenvironment{house@up}[2]% upright body: definition/notation/example/method/axiom/assumption/fact/remark
  {\par\addvspace{9pt}\refstepcounter{housethm}%
   \noindent{\bfseries\color{accent}#1~\thehousethm
     \if\relax\detokenize{#2}\relax\else\ (#2)\fi.}\hspace{0.4em}\ignorespaces}
  {\par\addvspace{9pt}}
\newenvironment{theorem}[1][]{\house@thm{Theorem}{#1}}{\endhouse@thm}
\newenvironment{lemma}[1][]{\house@thm{Lemma}{#1}}{\endhouse@thm}
\newenvironment{proposition}[1][]{\house@thm{Proposition}{#1}}{\endhouse@thm}
\newenvironment{corollary}[1][]{\house@thm{Corollary}{#1}}{\endhouse@thm}
\newenvironment{claim}[1][]{\house@thm{Claim}{#1}}{\endhouse@thm}
\newenvironment{definition}[1][]{\house@up{Definition}{#1}}{\endhouse@up}
\newenvironment{notation}[1][]{\house@up{Notation}{#1}}{\endhouse@up}
\newenvironment{example}[1][]{\house@up{Example}{#1}}{\endhouse@up}
\newenvironment{method}[1][]{\house@up{Method}{#1}}{\endhouse@up}
\newenvironment{axiom}[1][]{\house@up{Axiom}{#1}}{\endhouse@up}
\newenvironment{assumption}[1][]{\house@up{Assumption}{#1}}{\endhouse@up}
\newenvironment{fact}[1][]{\house@up{Fact}{#1}}{\endhouse@up}
\newenvironment{remark}[1][]{\house@up{Remark}{#1}}{\endhouse@up}
\makeatother
% \proofname exists only when amsthm/amsproof is loaded; define-or-override so
% this block stays self-contained in courses that load neither.
\providecommand{\proofname}{Proof}
\renewcommand{\proofname}{\color{accent}Proof}
            % geometry, colors, sections, ToC, heads
%     \input{theorems}              % house math environments
%     % ============================================================
% COMMON PREAMBLE — identical block for every course (do not diverge).
% Load order in each main.tex:
%     \documentclass[11pt,a4paper]{article}
%     \usepackage{lmodern}          % house serif font
%     \input{housestyle}            % geometry, colors, sections, ToC, heads
%     \input{theorems}              % house math environments
%     \input{common-preamble}       % THIS FILE — the shared package set
% Everything a course actually needs to compile lives here, so the package set
% is byte-identical across the collection. Packages a given course never calls
% are still loaded, on purpose, for coherence. What is NOT here (and why):
%   - lmodern .......... loaded in the main before housestyle (font must come first)
%   - geometry/xcolor/titlesec/fancyhdr/enumitem/caption/tocloft .. owned by housestyle
%   - lato (sans) ...... lids-only font choice; global load would restyle every heading
%   - \addbibresource .. the .bib path is course-specific; kept in each main
% ============================================================

% --- Encoding & language ---
\usepackage[T1]{fontenc}
\usepackage[utf8]{inputenc}
\usepackage[english]{babel}

% --- Math ---
\usepackage{amsmath, amssymb, amsthm, amsfonts, mathtools}

% --- Tables ---
\usepackage{booktabs}
\usepackage{tabularx}
\usepackage{array}
\usepackage{colortbl}
\usepackage{multirow}
\renewcommand{\arraystretch}{1.2}
\newcolumntype{L}{>{\raggedright\arraybackslash}X}
\newcolumntype{C}{>{\centering\arraybackslash}X}
\newcolumntype{R}{>{\raggedleft\arraybackslash}X}

% --- Graphics & figures ---
\usepackage{graphicx}
\usepackage{float}
\usepackage{wrapfig}
\usepackage{adjustbox}
\usepackage{subcaption}
\usepackage{transparent}

% --- Multi-column ---
\usepackage{multicol}

% --- TikZ ---
\usepackage{tikz}
\usetikzlibrary{shapes.geometric, arrows.meta, arrows, positioning, fit,
                backgrounds, calc, shapes}

% --- Listings (code / pseudocode) ---
\usepackage{listings}
\lstset{
  basicstyle=\ttfamily\footnotesize,
  backgroundcolor=\color{rowgray},
  frame=none,
  breaklines=true,
  showstringspaces=false,
  xleftmargin=4pt, xrightmargin=4pt,
  aboveskip=4pt, belowskip=4pt
}

% --- Algorithms ---
\usepackage{algorithm}
\usepackage{algpseudocode}

% --- Boxes ---
\usepackage[most]{tcolorbox}
\usepackage{mdframed}

% --- House box: one grammar for every framed callout across the collection.
%     Thin accent left band, faint accent tint, optional run-in title. Courses
%     map their own box names (notebox/theorybox/normref/...) onto this so a box
%     looks the same everywhere even when it says different things. ---
\tcbset{
  housebox/.style={
    enhanced, breakable,
    colback=accent!3, colframe=accent,
    colbacktitle=accent!3,          % title bar shares the box tint (no full band)
    boxrule=0pt, leftrule=1.5pt,
    arc=0pt, outer arc=0pt,
    left=8pt, right=8pt, top=5pt, bottom=5pt,
    before skip=8pt, after skip=8pt,
    fonttitle=\sffamily\bfseries\small,
    coltitle=accent,                % accent title text on the light tint
    titlerule=0pt,
    toptitle=1pt, bottomtitle=1pt,
  }
}
% Back-compat: old lds/lids boxes were mdframed and passed the title via the
% key `frametitle={...}`. Alias it to tcolorbox's `title` so the section bodies
% need no edits.
\tcbset{frametitle/.style={title={#1}}}
% One generic framed environment in the house grammar, named `calloutbox`. A
% course maps its own box names onto it (see each course.tex) instead of the
% common file claiming generic names like notebox/examplebox — those mean
% different things in different courses (run-in callouts in bpm/o4ds vs framed
% boxes in lds), so they must NOT be defined here.
\newtcolorbox{calloutbox}[1][]{housebox, #1}

% --- Symbols & misc ---
\usepackage{eurosym}
\usepackage{microtype}

% --- Cross-references (must stay near the end, before hyperref) ---
% cleveref is loaded by each main AFTER hyperref (load-order requirement).

% --- Bibliography: one backend for the whole collection (biblatex + biber,
%     numeric style). Courses with no .bib simply print nothing. \cite works. ---
\usepackage[backend=biber, style=numeric, sorting=none]{biblatex}
       % THIS FILE — the shared package set
% Everything a course actually needs to compile lives here, so the package set
% is byte-identical across the collection. Packages a given course never calls
% are still loaded, on purpose, for coherence. What is NOT here (and why):
%   - lmodern .......... loaded in the main before housestyle (font must come first)
%   - geometry/xcolor/titlesec/fancyhdr/enumitem/caption/tocloft .. owned by housestyle
%   - lato (sans) ...... lids-only font choice; global load would restyle every heading
%   - \addbibresource .. the .bib path is course-specific; kept in each main
% ============================================================

% --- Encoding & language ---
\usepackage[T1]{fontenc}
\usepackage[utf8]{inputenc}
\usepackage[english]{babel}

% --- Math ---
\usepackage{amsmath, amssymb, amsthm, amsfonts, mathtools}

% --- Tables ---
\usepackage{booktabs}
\usepackage{tabularx}
\usepackage{array}
\usepackage{colortbl}
\usepackage{multirow}
\renewcommand{\arraystretch}{1.2}
\newcolumntype{L}{>{\raggedright\arraybackslash}X}
\newcolumntype{C}{>{\centering\arraybackslash}X}
\newcolumntype{R}{>{\raggedleft\arraybackslash}X}

% --- Graphics & figures ---
\usepackage{graphicx}
\usepackage{float}
\usepackage{wrapfig}
\usepackage{adjustbox}
\usepackage{subcaption}
\usepackage{transparent}

% --- Multi-column ---
\usepackage{multicol}

% --- TikZ ---
\usepackage{tikz}
\usetikzlibrary{shapes.geometric, arrows.meta, arrows, positioning, fit,
                backgrounds, calc, shapes}

% --- Listings (code / pseudocode) ---
\usepackage{listings}
\lstset{
  basicstyle=\ttfamily\footnotesize,
  backgroundcolor=\color{rowgray},
  frame=none,
  breaklines=true,
  showstringspaces=false,
  xleftmargin=4pt, xrightmargin=4pt,
  aboveskip=4pt, belowskip=4pt
}

% --- Algorithms ---
\usepackage{algorithm}
\usepackage{algpseudocode}

% --- Boxes ---
\usepackage[most]{tcolorbox}
\usepackage{mdframed}

% --- House box: one grammar for every framed callout across the collection.
%     Thin accent left band, faint accent tint, optional run-in title. Courses
%     map their own box names (notebox/theorybox/normref/...) onto this so a box
%     looks the same everywhere even when it says different things. ---
\tcbset{
  housebox/.style={
    enhanced, breakable,
    colback=accent!3, colframe=accent,
    colbacktitle=accent!3,          % title bar shares the box tint (no full band)
    boxrule=0pt, leftrule=1.5pt,
    arc=0pt, outer arc=0pt,
    left=8pt, right=8pt, top=5pt, bottom=5pt,
    before skip=8pt, after skip=8pt,
    fonttitle=\sffamily\bfseries\small,
    coltitle=accent,                % accent title text on the light tint
    titlerule=0pt,
    toptitle=1pt, bottomtitle=1pt,
  }
}
% Back-compat: old lds/lids boxes were mdframed and passed the title via the
% key `frametitle={...}`. Alias it to tcolorbox's `title` so the section bodies
% need no edits.
\tcbset{frametitle/.style={title={#1}}}
% One generic framed environment in the house grammar, named `calloutbox`. A
% course maps its own box names onto it (see each course.tex) instead of the
% common file claiming generic names like notebox/examplebox — those mean
% different things in different courses (run-in callouts in bpm/o4ds vs framed
% boxes in lds), so they must NOT be defined here.
\newtcolorbox{calloutbox}[1][]{housebox, #1}

% --- Symbols & misc ---
\usepackage{eurosym}
\usepackage{microtype}

% --- Cross-references (must stay near the end, before hyperref) ---
% cleveref is loaded by each main AFTER hyperref (load-order requirement).

% --- Bibliography: one backend for the whole collection (biblatex + biber,
%     numeric style). Courses with no .bib simply print nothing. \cite works. ---
\usepackage[backend=biber, style=numeric, sorting=none]{biblatex}
       % THIS FILE — the shared package set
% Everything a course actually needs to compile lives here, so the package set
% is byte-identical across the collection. Packages a given course never calls
% are still loaded, on purpose, for coherence. What is NOT here (and why):
%   - lmodern .......... loaded in the main before housestyle (font must come first)
%   - geometry/xcolor/titlesec/fancyhdr/enumitem/caption/tocloft .. owned by housestyle
%   - lato (sans) ...... lids-only font choice; global load would restyle every heading
%   - \addbibresource .. the .bib path is course-specific; kept in each main
% ============================================================

% --- Encoding & language ---
\usepackage[T1]{fontenc}
\usepackage[utf8]{inputenc}
\usepackage[english]{babel}

% --- Math ---
\usepackage{amsmath, amssymb, amsthm, amsfonts, mathtools}

% --- Tables ---
\usepackage{booktabs}
\usepackage{tabularx}
\usepackage{array}
\usepackage{colortbl}
\usepackage{multirow}
\renewcommand{\arraystretch}{1.2}
\newcolumntype{L}{>{\raggedright\arraybackslash}X}
\newcolumntype{C}{>{\centering\arraybackslash}X}
\newcolumntype{R}{>{\raggedleft\arraybackslash}X}

% --- Graphics & figures ---
\usepackage{graphicx}
\usepackage{float}
\usepackage{wrapfig}
\usepackage{adjustbox}
\usepackage{subcaption}
\usepackage{transparent}

% --- Multi-column ---
\usepackage{multicol}

% --- TikZ ---
\usepackage{tikz}
\usetikzlibrary{shapes.geometric, arrows.meta, arrows, positioning, fit,
                backgrounds, calc, shapes}

% --- Listings (code / pseudocode) ---
\usepackage{listings}
\lstset{
  basicstyle=\ttfamily\footnotesize,
  backgroundcolor=\color{rowgray},
  frame=none,
  breaklines=true,
  showstringspaces=false,
  xleftmargin=4pt, xrightmargin=4pt,
  aboveskip=4pt, belowskip=4pt
}

% --- Algorithms ---
\usepackage{algorithm}
\usepackage{algpseudocode}

% --- Boxes ---
\usepackage[most]{tcolorbox}
\usepackage{mdframed}

% --- House box: one grammar for every framed callout across the collection.
%     Thin accent left band, faint accent tint, optional run-in title. Courses
%     map their own box names (notebox/theorybox/normref/...) onto this so a box
%     looks the same everywhere even when it says different things. ---
\tcbset{
  housebox/.style={
    enhanced, breakable,
    colback=accent!3, colframe=accent,
    colbacktitle=accent!3,          % title bar shares the box tint (no full band)
    boxrule=0pt, leftrule=1.5pt,
    arc=0pt, outer arc=0pt,
    left=8pt, right=8pt, top=5pt, bottom=5pt,
    before skip=8pt, after skip=8pt,
    fonttitle=\sffamily\bfseries\small,
    coltitle=accent,                % accent title text on the light tint
    titlerule=0pt,
    toptitle=1pt, bottomtitle=1pt,
  }
}
% Back-compat: old lds/lids boxes were mdframed and passed the title via the
% key `frametitle={...}`. Alias it to tcolorbox's `title` so the section bodies
% need no edits.
\tcbset{frametitle/.style={title={#1}}}
% One generic framed environment in the house grammar, named `calloutbox`. A
% course maps its own box names onto it (see each course.tex) instead of the
% common file claiming generic names like notebox/examplebox — those mean
% different things in different courses (run-in callouts in bpm/o4ds vs framed
% boxes in lds), so they must NOT be defined here.
\newtcolorbox{calloutbox}[1][]{housebox, #1}

% --- Symbols & misc ---
\usepackage{eurosym}
\usepackage{microtype}

% --- Cross-references (must stay near the end, before hyperref) ---
% cleveref is loaded by each main AFTER hyperref (load-order requirement).

% --- Bibliography: one backend for the whole collection (biblatex + biber,
%     numeric style). Courses with no .bib simply print nothing. \cite works. ---
\usepackage[backend=biber, style=numeric, sorting=none]{biblatex}
       % THIS FILE — the shared package set
% Everything a course actually needs to compile lives here, so the package set
% is byte-identical across the collection. Packages a given course never calls
% are still loaded, on purpose, for coherence. What is NOT here (and why):
%   - lmodern .......... loaded in the main before housestyle (font must come first)
%   - geometry/xcolor/titlesec/fancyhdr/enumitem/caption/tocloft .. owned by housestyle
%   - lato (sans) ...... lids-only font choice; global load would restyle every heading
%   - \addbibresource .. the .bib path is course-specific; kept in each main
% ============================================================

% --- Encoding & language ---
\usepackage[T1]{fontenc}
\usepackage[utf8]{inputenc}
\usepackage[english]{babel}

% --- Math ---
\usepackage{amsmath, amssymb, amsthm, amsfonts, mathtools}

% --- Tables ---
\usepackage{booktabs}
\usepackage{tabularx}
\usepackage{array}
\usepackage{colortbl}
\usepackage{multirow}
\renewcommand{\arraystretch}{1.2}
\newcolumntype{L}{>{\raggedright\arraybackslash}X}
\newcolumntype{C}{>{\centering\arraybackslash}X}
\newcolumntype{R}{>{\raggedleft\arraybackslash}X}

% --- Graphics & figures ---
\usepackage{graphicx}
\usepackage{float}
\usepackage{wrapfig}
\usepackage{adjustbox}
\usepackage{subcaption}
\usepackage{transparent}

% --- Multi-column ---
\usepackage{multicol}

% --- TikZ ---
\usepackage{tikz}
\usetikzlibrary{shapes.geometric, arrows.meta, arrows, positioning, fit,
                backgrounds, calc, shapes}

% --- Listings (code / pseudocode) ---
\usepackage{listings}
\lstset{
  basicstyle=\ttfamily\footnotesize,
  backgroundcolor=\color{rowgray},
  frame=none,
  breaklines=true,
  showstringspaces=false,
  xleftmargin=4pt, xrightmargin=4pt,
  aboveskip=4pt, belowskip=4pt
}

% --- Algorithms ---
\usepackage{algorithm}
\usepackage{algpseudocode}

% --- Boxes ---
\usepackage[most]{tcolorbox}
\usepackage{mdframed}

% --- House box: one grammar for every framed callout across the collection.
%     Thin accent left band, faint accent tint, optional run-in title. Courses
%     map their own box names (notebox/theorybox/normref/...) onto this so a box
%     looks the same everywhere even when it says different things. ---
\tcbset{
  housebox/.style={
    enhanced, breakable,
    colback=accent!3, colframe=accent,
    colbacktitle=accent!3,          % title bar shares the box tint (no full band)
    boxrule=0pt, leftrule=1.5pt,
    arc=0pt, outer arc=0pt,
    left=8pt, right=8pt, top=5pt, bottom=5pt,
    before skip=8pt, after skip=8pt,
    fonttitle=\sffamily\bfseries\small,
    coltitle=accent,                % accent title text on the light tint
    titlerule=0pt,
    toptitle=1pt, bottomtitle=1pt,
  }
}
% Back-compat: old lds/lids boxes were mdframed and passed the title via the
% key `frametitle={...}`. Alias it to tcolorbox's `title` so the section bodies
% need no edits.
\tcbset{frametitle/.style={title={#1}}}
% One generic framed environment in the house grammar, named `calloutbox`. A
% course maps its own box names onto it (see each course.tex) instead of the
% common file claiming generic names like notebox/examplebox — those mean
% different things in different courses (run-in callouts in bpm/o4ds vs framed
% boxes in lds), so they must NOT be defined here.
\newtcolorbox{calloutbox}[1][]{housebox, #1}

% --- Symbols & misc ---
\usepackage{eurosym}
\usepackage{microtype}

% --- Cross-references (must stay near the end, before hyperref) ---
% cleveref is loaded by each main AFTER hyperref (load-order requirement).

% --- Bibliography: one backend for the whole collection (biblatex + biber,
%     numeric style). Courses with no .bib simply print nothing. \cite works. ---
\usepackage[backend=biber, style=numeric, sorting=none]{biblatex}
